% Part: lambda-calculus
% Chapter: syntax
% Section: free-variables

\documentclass[../../../include/open-logic-section]{subfiles}

\begin{document}

\olfileid{lam}{syn}{fv}
\olsection{متغیرهای آزاد}

حساب لامبدا دربارهٔ توابع است و انتزاع لامبدایی شیوهٔ پدیدآمدن توابع
است. به‌طور شهودی، $\lambd[x][M]$ تابعی است که وقتی آرگومان تابع به~$x$
اختصاص داده شود، ارزش‌هایش با~$M$ تعیین می‌شوند. اما هر رخداد~$x$ در~$M$
مرتبط نیست: اگر~$M$ انتزاع دیگری چون~$\lambd[x][N]$ داشته باشد، رخدادهای
$x$ در~$N$ به~$\lambd[x][N]$ مربوط‌اند، نه به~$\lambd[x][M]$. پس انتزاع
لامبدای~$\lambd[x]$ درون~$\lambd[x][M]$ رخدادهایی از~$x$ در~$M$ را که
از پیش با انتزاع لامبدای دیگری مقید نشده‌اند، یعنی رخدادهای \emph{آزاد}
$x$ در~$M$، \emph{مقید می‌کند}.

\begin{defn}[دامنه]
اگر~$\lambd[x][M]$ درون ترم~$N$ رخ دهد، رخداد متناظر~$M$، \emph{دامنهٔ}
$\lambd[x]$ است.
\end{defn}


\begin{defn}[رخداد آزاد و مقید]
  رخداد متغیر~$x$ در ترم~$M$ \emph{آزاد} است اگر در دامنهٔ یک
  $\lambd[x]$ نباشد، وگرنه \emph{مقید} است. رخداد متغیر~$x$ در
  $\lambd[x][M]$ با~$\lambd[x]$ آغازین مقید است اگر و تنها اگر رخداد~$x$
  در~$M$ آزاد باشد.
\end{defn}

\begin{ex}
  در~$\lambd[x][x y]$، هر دو~$x$ و~$y$ در دامنهٔ~$\lambd[x]$ هستند، پس
  $x$ با~$\lambd[x]$ مقید است. چون~$y$ در دامنهٔ هیچ~$\lambd[y]$ای
  نیست، آزاد است. در~$\lambd[x][x x]$ هر دو رخداد~$x$ با~$\lambd[x]$
  مقیدند، زیرا هر دو در~$xx$ آزادند. در~$((\lambd[x][xx])x)$، رخداد آخر
  $x$ آزاد است، زیرا در دامنهٔ هیچ~$\lambd[x]$ای نیست. در
  $\lambd[x][(\lambd[x][x])x]$، دامنهٔ~$\lambd[x]$ نخست
  $(\lambd[x][x])x$ و دامنهٔ~$\lambd[x]$ دوم رخداد یکی‌مانده‌به‌آخر~$x$
  است. در~$(\lambd[x][x])x$، رخداد آخر~$x$ آزاد و رخداد
  یکی‌مانده‌به‌آخر مقید است. پس رخداد یکی‌مانده‌به‌آخر~$x$ در
  $\lambd[x][(\lambd[x][x])x]$ با~$\lambd[x]$ دوم و رخداد آخر با
  $\lambd[x]$ نخست مقید است.
\end{ex}

برای ترم~$P$ می‌توانیم همهٔ رخدادهای متغیر در آن را بررسی کنیم و مجموعهٔ
متغیرهای آزاد را به دست آوریم. این مجموعه با~$\FV{P}$ نشان داده می‌شود
و تعریف طبیعی زیر را دارد:

\begin{defn}[متغیرهای آزاد یک ترم] \ollabel{def:fv}
  مجموعهٔ \emph{متغیرهای آزاد} یک ترم به‌صورت استقرایی چنین تعریف می‌شود:
  \begin{enumerate}
    \item $\FV{x} = \{x\}$ \ollabel{def:fv1}
    \item $\FV{\lambd[x][N]} = \FV{N} \setminus \{x\}$    \ollabel{def:fv2}
    \item $\FV{PQ} = \FV{P} \cup \FV{Q}$ \ollabel{def:fv3}
  \end{enumerate}
\end{defn}

\begin{prob}
  \begin{enumerate}
  \item دامنه‌های~$\lambd[g]$ و دو~$\lambd[x]$ را در ترم زیر مشخص کنید:
    $\lambd[g][(\lambd[x][g (x x)]) \lambd[x][g (x x)]]$.
  \item آیا همهٔ رخدادهای متغیرها در
    $\lambd[g][(\lambd[x][g (x x)]) \lambd[x][g (x x)]]$ مقیدند؟ هر
    یک به‌ترتیب با کدام انتزاع مقید شده‌اند؟
    \item مقدار~$\FV{\lambd[x][(\lambd[y][(\lambd[z][x y]) z]) y]}$
    را به دست آورید.
  \end{enumerate}
\end{prob}

\begin{explain}
متغیر آزاد مانند ارجاعی به جهان بیرون (\emph{محیط}) است و ترمی که متغیر
آزاد دارد را می‌توان ترمی با مشخصات جزئی دانست، زیرا رفتارش به چگونگی
تنظیم محیط بستگی دارد. برای مثال، در ترم~$\lambd[x][f x]$ که آرگومان~$x$
را می‌پذیرد و حاصل اعمال~$f$ بر آن آرگومان را برمی‌گرداند، متغیر~$f$ آزاد است. ارزش
این ترم به محیطی که در آن قرار دارد و به‌ویژه به ارزش~$f$ در آن محیط
بستگی دارد.

اگر انتزاع را بر این ترم اعمال کنیم، $\lambd[f][\lambd[x][f
    x]]$ را به دست می‌آوریم. این ترم دیگر به متغیر محیطی~$f$ وابسته نیست،
زیرا اکنون تابعی را نشان می‌دهد که دو آرگومان می‌پذیرد و نتیجهٔ اعمال
اولی بر دومی را برمی‌گرداند. تغییر~$f$ در محیط اثری بر رفتار این ترم
ندارد، زیرا ترم تنها از چیزی استفاده می‌کند که به‌عنوان آرگومان به آن
داده شده است، نه از ارزش~$f$ در محیط.
\end{explain}

\begin{defn}[ترم بسته، ترکیب‌گر]
  ترمی که متغیر آزاد ندارد \emph{ترم بسته} یا \emph{ترکیب‌گر} نامیده
  می‌شود.
\end{defn}

\begin{lem}\ollabel{lem:fv}
  \begin{enumerate}
    \item \ollabel{lem:fv-abs} اگر~$y \neq x$، آنگاه~$y \in
      \FV{\lambd[x][N]}$ اگر و تنها اگر~$y \in \FV{N}$.
    \item \ollabel{lem:fv-app} رابطهٔ~$y \in \FV{PQ}$ برقرار است اگر
      و تنها اگر~$y \in \FV{P}$ یا~$y \in \FV{Q}$.
    \end{enumerate}
\end{lem}

\begin{proof}
  تمرین.
\end{proof}

\begin{prob}
  \olref[lam][syn][fv]{lem:fv} را ثابت کنید.
\end{prob}

\end{document}
